
\section{Introduction}
 
Large language models (LLMs), such as GPT \cite{DBLP:journals/corr/abs-2303-08774}, Claude \cite{anthropic2023claude} and Llama \cite{DBLP:journals/corr/abs-2307-09288}, have shown remarkable performance on a wide range of natural language tasks. However, when the target applications are in highly specialized domains, such as Medicine, Finance, Law, etc., general-purpose LLMs often struggle with specialized terminologies and domain‑specific reasoning needs. Consequently, researchers have begun to build domain‑specific variants by specifically pre‑training and fine-tuning on corpora from the corresponding target fields \cite{singhal2023large,Singhal:2025aa,xie-etal-2024-efficient,DBLP:conf/nips/ColomboPBMHMCCD24}.

% In Medicine, Med-PaLM/Med-PaLM-2 \cite{singhal2023large,Singhal:2025aa} was instruction-tuned from PaLM/Flan-PaLM on medical QA datasets and clinical guidelines, reaching exam-level performance on MedQA-style tasks. In Finance, FinPythia-6.9B \cite{xie-etal-2024-efficient} showed how domain-adaptive continual pre-training on financial texts led to superior performance on financial NLP tasks over the base model. In Legal, several recent works have highlighted the need for dedicated legal LLMs. \cite{jayakumar-etal-2023-large,DBLP:journals/corr/abs-2304-12202} showed that general‑purpose models lag behind smaller fine-tuned legal models, while domain-specific models like SAUL \cite{NEURIPS2024_ea3f85a3} achieved improvements on various legal benchmarks.

While training and continuous fine-tuning on raw text capture the generation of domain language, they are insufficient to encode the structure of domain knowledge, especially on domain-specific reasoning tasks. Recent reasoning models work best for areas, such as mathematics and programming, where verifiable rewards enable a straightforward approach to reinforcement learning; however, areas like Medicine or Law require alternative forms of reasoning, such as argumentative reasoning or inference in the presence of high and inherent uncertainty, by using concepts of the corresponding domains, such as symptoms, events, and legal rules.

More importantly, these areas require deep understanding of the relations between abstract concepts. Financial market-event analysis adopts a cause-effect-impact-conclusion chain; clinical reasoning follows a symptom-diagnosis-treatment-outcome flow; in Legal, judicial reasoning employs the fact-rule-conclusion process. Encoding these reasoning patterns in a structured form offers a principled way to augment LLMs' inductive bias and to reduce hallucinations or misattributions.

In this paper, we focus on the legal domain where even when models are trained on vast corpora of case law, statutes, regulations, and secondary literature, they may produce misattributions, hallucinations, and logically incoherent arguments \cite{doi.org/10.1111/jels.12413}, and specific tasks are often better served by smaller specialized models \cite{DBLP:journals/corr/abs-2304-12202,jayakumar-etal-2023-large,DBLP:conf/cikm/ShuZLDDZ24}. The root of this problem lies not in the quantity of text for LLM training but in the structure of legal knowledge: legal prose is dense, highly contextual, and governed by rules that change over time, calling for a clearer textual representation that makes the element of the reasoning process more accessible.

In order to address this gap, we start with the formal modeling of legal concepts, which then leads to generating textual training data for LLM post-training. Our approach explicitly encodes the reasoning patterns found in judicial opinions in order to facilitate LLM training. The cornerstone of legal reasoning is the \textbf{IRAC} framework, \textit{\textbf{I}ssue}, \textit{\textbf{R}ule}, \textit{\textbf{A}nalysis} and \textit{\textbf{C}onclusion}, which is taught in law schools and encapsulates the typical flow of legal argumentation \cite{brand1976legalwriting}. By constructing an \textbf{IRAC Knowledge Graph} (IRAC KG) that captures the relationships between issues, governing rules, factual applications, and conclusions, we aim to provide a structured substrate from which an LLM can learn to reason more faithfully. 

Our work equips LLMs with a structured substrate for legal reasoning, but with two distinct features. First, instead of focusing on creating legal KGs with specific entities, such as Person, Organization, etc., our goal is to model the key legal concepts by following the IRAC framework. Second, rather than treating the graph merely as retrieval scaffolding, we post-train LLMs directly on graph-derived supervision signals, and show that this yields measurable gains on diverse legal tasks. Our contributions are threefold:
\begin{enumerate}
    \item \textbf{IRAC KG construction.} We extract IRAC components from a large corpus of 12K legal case opinions, and encode them as nodes and typed edges in a graph, preserving the causal dependencies that underlie legal reasoning.
    
    \item \textbf{IRAC KG-based LLM post-training.} By utilizing the KG as a knowledge source, we post-train mid to large LLMs (Qwen3-30B-A3B-Instruct-2507, Llama-3\_3-Nemotron-Super-49B-v1\_5 and Llama-3.1-70B-Instruct) via both Supervised Fine-Tuning (SFT) and Direct Preference Optimization (DPO) \cite{DBLP:conf/nips/RafailovSMMEF23} in order to improve their ability to generate coherent legal arguments.
    
    \item \textbf{Extensive evaluation.} We assess the trained models on five legal benchmarks: LexGLUE \cite{DBLP:journals/corr/abs-2304-12202}, LegalBench \cite{DBLP:conf/nips/GuhaNHRCKCPWRZT23}, COLIEE \cite{DBLP:journals/rss/RabeloGKKYS22}, SuperGPQA (Law) \cite{DBLP:journals/corr/abs-2502-14739} and $\delta$-Stance \cite{DBLP:conf/acl/GuptaR025}, demonstrating measurable gains over baseline and SOTA models.
\end{enumerate}

%We organize the rest of the paper as follows. In Section \ref{section:related_work}, we discuss related work to KG, LLM and their domain-specific variants. We then formally present our IRAC KG in Section \ref{section:overview_and_kg_design} and model training details in Section \ref{section:data_generation}. We present and discuss our evaluation results in Section \ref{section:evaluation} and conclude in \ref{section:conclusion}.
